We identified three robust PD progression subtypes through unsupervised trajectory clustering of 536 patients followed over 5 years, demonstrating that data-driven subtyping captures clinically and biologically meaningful heterogeneity. These subtypes—rapid motor decline (22\%), moderate progression (54\%), and stable motor with cognitive decline (24\%)—exhibited distinct baseline predictors, biomarker profiles, and natural histories. Critically, subtype-based trial enrichment reduced required sample sizes by 38-69\%, providing actionable strategies for precision medicine and adaptive trial designs.

Our findings align with and extend prior PD subtyping efforts. Fereshtehnejad et al.\cite{fereshtehnejad2017subtypes} identified cross-sectional motor phenotypes with prognostic implications, while Lawton et al.\cite{lawton2018developing} reported 25\% rapid progressors using threshold-based categorization. Our unsupervised approach validates these proportions (22\% rapid) while revealing a distinct cognitive-first phenotype (Subtype 3, 24\%) not captured by motor-centric classifications. The cognitive decline subtype, characterized by minimal motor progression but rapid MoCA decline (-1.8 points/year), may represent Lewy body pathology with earlier cortical spread\cite{braak2003staging}, warranting targeted cognitive interventions.

The partial predictability of subtypes from baseline features (68\% accuracy, AUC 0.74) indicates that progression heterogeneity arises from both baseline characteristics and stochastic disease processes. Top predictors—baseline UPDRS-III, age, MoCA, sex, DaTscan SBR—are routinely collected, enabling prospective stratification without specialized assays. The 83\% accuracy for binary rapid vs. non-rapid classification suffices for trial enrichment, as demonstrated by simulations. Notably, biomarkers improved prediction over clinical-only models (74\% vs. 68\% AUC), justifying their inclusion in prognostic algorithms.

Biomarker associations validate biological distinctiveness. Rapid progressors exhibited more severe dopaminergic deficits (lowest DaTscan SBR), reduced CSF $\alpha$-synuclein (potentially indicating greater aggregation\cite{kang2016biofind}), and elevated GBA mutation carriage (14\% vs. 6-8\%)—consistent with GBA-associated accelerated progression\cite{alcalay2012cognitive}. These associations suggest subtype membership reflects underlying pathophysiology rather than measurement artifacts, supporting etiological heterogeneity.

The trial enrichment implications are substantial. A 38\% sample size reduction (moderate enrichment) translates to $\sim$170 fewer patients for typical trials ($N=424\to 254$ per arm), reducing costs by an estimated \$15-20 million and accelerating enrollment timelines. Rapid-only enrichment achieves even greater efficiency (69\% reduction) but may limit generalizability and recruitment feasibility. We therefore recommend moderate enrichment or subtype-stratified designs, balancing efficiency with external validity. Importantly, enrichment does not compromise statistical validity—stratified randomization ensures balanced group allocation while leveraging prognostic information to boost power\cite{kernan1999stratified}.

Our approach has limitations. First, \ppmi{} is a research cohort with rigorous inclusion criteria, potentially limiting generalizability to community populations with greater comorbidity. External validation in BioFIND confirmed similar subtype proportions (20\% rapid), but broader validation in diverse cohorts is needed. Second, trajectory clustering assumes piecewise-linear progression, not capturing nonlinear patterns (e.g., sudden declines). More flexible models (e.g., Gaussian process latent variable models\cite{fortuin2020gp}) could address this. Third, 5-year follow-up may inadequately capture late-stage heterogeneity. Longer follow-up (10-15 years) may reveal additional subtypes or subtype transitions\cite{Post2018}. Fourth, we lacked multimodal imaging (e.g., MRI cortical thickness) which may further refine subtypes\cite{zeighami2019network}.

Clinical translation requires validated prediction algorithms deployable at point-of-care. We are developing a web-based calculator (\url{pdsubtypes.org}, in progress) providing individualized subtype probabilities and progression estimates from baseline features. Prospective validation studies (e.g., embedding prediction algorithms in ongoing trials) are essential before clinical adoption. Additionally, subtype-specific interventions warrant investigation: rapid progressors may benefit from aggressive dopaminergic therapy escalation or neuroprotective trials, while cognitive-first patients may require earlier cognitive training or dementia therapies.

The methodological framework—VAE trajectory modeling, latent time alignment, silhouette-validated clustering—generalizes beyond PD to other neurodegenerative diseases exhibiting longitudinal heterogeneity (Alzheimer's disease\cite{young2018longitudinal}, multiple sclerosis\cite{galea2018machine}, ALS\cite{westeneng2018prognosis}). The trial enrichment principles apply broadly: identifying fast-progressing subgroups amplifies treatment signals, improving statistical efficiency. However, enrichment strategies must be weighed against potential selection bias, generalizability limitations, and regulatory considerations regarding label indications.

Future work should integrate genetic risk scores (polygenic risk scores from recent GWAS\cite{nalls2019identification}), multimodal imaging (MRI, tau-PET\cite{passamonti2017tau}), and multi-omics data (metabolomics, proteomics\cite{maass2020multimodal}) to refine subtype prediction and elucidate biological mechanisms. Longitudinal multi-omic profiling may identify subtype-specific pathways amenable to targeted therapies. Additionally, survival-based subtyping (clustering time-to-event outcomes like time to dementia or H\&Y stage 5) could directly inform prognosis. Machine learning advances—particularly deep learning for time series\cite{fortuin2020gp} and causal inference methods\cite{schulam2019can}—promise further improvements.

In conclusion, we demonstrate that data-driven longitudinal subtyping identifies robust PD progression phenotypes with distinct biological underpinnings, baseline predictability, and substantial implications for precision prognostication and clinical trial efficiency. Subtype-based trial enrichment offers a pragmatic strategy to accelerate therapeutic development, reducing costs and timelines while maintaining statistical rigor. These findings establish a roadmap for implementing precision medicine in PD, with applicability to broader neurodegenerative disease research.
